% GENERATED FILE. Edit canonical lessons, then run npm run generate:books.
% canonical-source-hash: fnv1a64:358d1bc4214258e0
% canonical-lessons: TE-C388-sandeham, TE-C388-sandehincu, TE-C388-visvasincu, TE-C388-niscayam, TE-C388-anisciti

\chapter{Doubt and Certainty}
\label{ch:te-doubt-and-certainty}
\hlchaptermodality{388}

\begin{chapteropening}
\textbf{By the end of this chapter:} \emph{I can talk about doubt, trust and certainty.}

Complete the last lesson of chapter 388: I can talk about doubt, trust and certainty.
\end{chapteropening}

\section[\texorpdfstring{\te{సందేహం} (sandēhaṁ)}{సందేహం (sandēhaṁ)}]{\te{సందేహం} (sandēhaṁ) — a misgiving; a doubt}

\label{lesson:TE-C388-sandeham}



\begin{quote}
Before the new one: say the Telugu for to comment, then the Telugu for to find fault with.
\end{quote}

\subsection*{You'll want to know: \te{సందేహం}}
\textbf{\te{సందేహం}} — \emph{sandēhaṁ} — \textquotedblleft{}a misgiving; a doubt\textquotedblright{}.

\begin{quote}
\te{నాకు} \te{ఒక} \te{చిన్న} \te{సందేహం} \te{ఉంది}.
\end{quote}

\begin{quote}
\emph{nāku oka cinna sandēhaṁ undi.}
\end{quote}

\begin{quote}
I have a small doubt.
\end{quote}

Plural: \te{సందేహాలు}, \emph{sandēhālu}.

In class it is the polite word for a question you are unsure about. \te{అనుమానం}, which you know, leans to suspicion.

\begin{cousinweb}
\textbf{\te{సందేహం}} (\emph{sandēhaṁ}) is borrowed from Sanskrit \emph{saṁdeha}, \textquotedblleft{}doubt, uncertainty\textquotedblright{} (Monier-Williams).
\end{cousinweb}

\begin{grammarlens}[title={What you've built}]
One more word for giving reasons and opinions.
\end{grammarlens}

\subsection*{Your turn}
\begin{itemize}
\raggedright
  \item \emph{Say it:} \emph{sandēhaṁ}
  \item \emph{Say it:} \emph{sandēhaṁ}, once more
  \item \emph{Say it:} \emph{nāku oka cinna sandēhaṁ undi.} — I have a small doubt.
  \item \emph{Recall:} say the Telugu for to comment, then the Telugu for to find fault with, then say \emph{sandēhaṁ} again
\end{itemize}

\begin{tcolorbox}[breakable,colback=teal!4,colframe=teal!35!black,title={Before you move on}]
What does \te{సందేహం} mean? (\textquotedblleft{}A misgiving; a doubt\textquotedblright{}.) Say it once more.
\end{tcolorbox}
\section[\texorpdfstring{\te{సందేహించు} (sandēhiñcu)}{సందేహించు (sandēhiñcu)}]{\te{సందేహించు} (sandēhiñcu) — to have misgivings}

\label{lesson:TE-C388-sandehincu}



\begin{quote}
Before the new one: say the Telugu for to find fault with, then the Telugu for a misgiving.
\end{quote}

\subsection*{You'll want to know: \te{సందేహించు}}
\textbf{\te{సందేహించు}} — \emph{sandēhiñcu} — \textquotedblleft{}to have misgivings\textquotedblright{}.

\begin{quote}
\te{అతను} \te{నిజం} \te{చెప్పాడా} \te{అని} \te{నేను} \te{సందేహించాను}.
\end{quote}

\begin{quote}
\emph{atanu nijaṁ ceppāḍā ani nēnu sandēhiñcānu.}
\end{quote}

\begin{quote}
I doubted whether he had told the truth.
\end{quote}

Past (I …): \te{సందేహించాను}, \emph{sandēhiñcānu}.

What you doubt is often a question with \te{అని}: \te{చెప్పాడా} \te{అని} \te{సందేహించు}.

\begin{cousinweb}
\textbf{\te{సందేహించు}} (\emph{sandēhiñcu}) is built on Sanskrit \emph{saṁdeha}, \textquotedblleft{}doubt\textquotedblright{} (Monier-Williams), with the Telugu verb ending -\te{ఇంచు}; the noun is \te{సందేహం}.
\end{cousinweb}

\begin{grammarlens}[title={What you've built}]
One more word for giving reasons and opinions.
\end{grammarlens}

\subsection*{Your turn}
\begin{itemize}
\raggedright
  \item \emph{Say it:} \emph{sandēhiñcu}
  \item \emph{Say it:} \emph{sandēhiñcu}, once more
  \item \emph{Say it:} \emph{atanu nijaṁ ceppāḍā ani nēnu sandēhiñcānu.} — I doubted whether he had told the truth.
  \item \emph{Recall:} say the Telugu for to find fault with, then the Telugu for a misgiving, then say \emph{sandēhiñcu} again
\end{itemize}

\begin{tcolorbox}[breakable,colback=teal!4,colframe=teal!35!black,title={Before you move on}]
What does \te{సందేహించు} mean? (\textquotedblleft{}To have misgivings\textquotedblright{}.) Say it once more.
\end{tcolorbox}
\section[\texorpdfstring{\te{విశ్వసించు} (viśvasiñcu)}{విశ్వసించు (viśvasiñcu)}]{\te{విశ్వసించు} (viśvasiñcu) — to believe in; to have faith in}

\label{lesson:TE-C388-visvasincu}



\begin{quote}
Before the new one: say the Telugu for a misgiving, then the Telugu for to have misgivings.
\end{quote}

\subsection*{You'll want to know: \te{విశ్వసించు}}
\textbf{\te{విశ్వసించు}} — \emph{viśvasiñcu} — \textquotedblleft{}to believe in; to have faith in\textquotedblright{}.

\begin{quote}
\te{నేను} \te{నా} \te{స్నేహితులను} \te{పూర్తిగా} \te{విశ్వసిస్తాను}.
\end{quote}

\begin{quote}
\emph{nēnu nā snēhitulanu pūrtigā viśvasistānu.}
\end{quote}

\begin{quote}
I trust my friends completely.
\end{quote}

Past (I …): \te{విశ్వసించాను}, \emph{viśvasiñcānu}.

Formal beside \te{నమ్ము}. The noun is \te{విశ్వాసం}, trust.

\begin{cousinweb}
\textbf{\te{విశ్వసించు}} (\emph{viśvasiñcu}) is built on Sanskrit \emph{viśvāsa}, \textquotedblleft{}trust, confidence\textquotedblright{} (CDIAL 11966), with the Telugu verb ending -\te{ఇంచు}.
\end{cousinweb}

\begin{grammarlens}[title={What you've built}]
One more word for giving reasons and opinions.
\end{grammarlens}

\subsection*{Your turn}
\begin{itemize}
\raggedright
  \item \emph{Say it:} \emph{viśvasiñcu}
  \item \emph{Say it:} \emph{viśvasiñcu}, once more
  \item \emph{Say it:} \emph{nēnu nā snēhitulanu pūrtigā viśvasistānu.} — I trust my friends completely.
  \item \emph{Recall:} say the Telugu for a misgiving, then the Telugu for to have misgivings, then say \emph{viśvasiñcu} again
\end{itemize}

\begin{tcolorbox}[breakable,colback=teal!4,colframe=teal!35!black,title={Before you move on}]
What does \te{విశ్వసించు} mean? (\textquotedblleft{}To believe in; to have faith in\textquotedblright{}.) Say it once more.
\end{tcolorbox}
\section[\texorpdfstring{\te{నిశ్చయం} (niścayaṁ)}{నిశ్చయం (niścayaṁ)}]{\te{నిశ్చయం} (niścayaṁ) — certainty; a firm resolve}

\label{lesson:TE-C388-niscayam}



\begin{quote}
Before the new one: say the Telugu for to have misgivings, then the Telugu for to believe in.
\end{quote}

\subsection*{You'll want to know: \te{నిశ్చయం}}
\textbf{\te{నిశ్చయం}} — \emph{niścayaṁ} — \textquotedblleft{}certainty; a firm resolve\textquotedblright{}.

\begin{quote}
\te{అతను} \te{వస్తాడని} \te{నాకు} \te{నిశ్చయం} \te{లేదు}.
\end{quote}

\begin{quote}
\emph{atanu vastāḍani nāku niścayaṁ lēdu.}
\end{quote}

\begin{quote}
I am not certain that he will come.
\end{quote}

A family word too: \te{పెళ్ళి} \te{నిశ్చయం}, the fixing of a wedding.

\begin{cousinweb}
\textbf{\te{నిశ్చయం}} (\emph{niścayaṁ}) is borrowed from Sanskrit \emph{niścaya}, \textquotedblleft{}inquiry, conviction\textquotedblright{} (CDIAL 7439), which Monier-Williams also glosses \textquotedblleft{}certainty; a fixed resolve\textquotedblright{}.
\end{cousinweb}

\begin{grammarlens}[title={What you've built}]
One more word for giving reasons and opinions.
\end{grammarlens}

\subsection*{Your turn}
\begin{itemize}
\raggedright
  \item \emph{Say it:} \emph{niścayaṁ}
  \item \emph{Say it:} \emph{niścayaṁ}, once more
  \item \emph{Say it:} \emph{atanu vastāḍani nāku niścayaṁ lēdu.} — I am not certain that he will come.
  \item \emph{Recall:} say the Telugu for to have misgivings, then the Telugu for to believe in, then say \emph{niścayaṁ} again
\end{itemize}

\begin{tcolorbox}[breakable,colback=teal!4,colframe=teal!35!black,title={Before you move on}]
What does \te{నిశ్చయం} mean? (\textquotedblleft{}Certainty; a firm resolve\textquotedblright{}.) Say it once more.
\end{tcolorbox}
\section[\texorpdfstring{\te{అనిశ్చితి} (aniściti)}{అనిశ్చితి (aniściti)}]{\te{అనిశ్చితి} (aniściti) — uncertainty}

\label{lesson:TE-C388-anisciti}



\begin{quote}
Before the new one: say the Telugu for to believe in, then the Telugu for certainty.
\end{quote}

\subsection*{You'll want to know: \te{అనిశ్చితి}}
\textbf{\te{అనిశ్చితి}} — \emph{aniściti} — \textquotedblleft{}uncertainty\textquotedblright{}.

\begin{quote}
\te{ఉద్యోగాల} \te{విషయంలో} \te{చాలా} \te{అనిశ్చితి} \te{ఉంది}.
\end{quote}

\begin{quote}
\emph{udyōgāla viṣayaṁlō cālā aniściti undi.}
\end{quote}

\begin{quote}
There is a lot of uncertainty about jobs.
\end{quote}

\te{అ}- means not: \te{అనిశ్చితి} is the lack of \te{నిశ్చయం}.

\begin{cousinweb}
\textbf{\te{అనిశ్చితి}} (\emph{aniściti}) is built from Sanskrit \emph{a-}, \textquotedblleft{}not\textquotedblright{}, and \emph{niściti}, \textquotedblleft{}ascertainment, settling\textquotedblright{} (Monier-Williams), a word of the same family as \te{నిశ్చయం}.
\end{cousinweb}

\begin{grammarlens}[title={What you've built}]
One more word for giving reasons and opinions.
\end{grammarlens}

\subsection*{Your turn}
\begin{itemize}
\raggedright
  \item \emph{Say it:} \emph{aniściti}
  \item \emph{Say it:} \emph{aniściti}, once more
  \item \emph{Say it:} \emph{udyōgāla viṣayaṁlō cālā aniściti undi.} — There is a lot of uncertainty about jobs.
  \item \emph{Recall:} say the Telugu for to believe in, then the Telugu for certainty, then say \emph{aniściti} again
\end{itemize}

\begin{tcolorbox}[breakable,colback=teal!4,colframe=teal!35!black,title={Before you move on}]
What does \te{అనిశ్చితి} mean? (\textquotedblleft{}Uncertainty\textquotedblright{}.) Say it once more.
\end{tcolorbox}
